\section{A determinant section of the rank-three deformation space}

Let \(O,F,C_0\) be the original coefficient DVR, frame and parameter
with \(F.\mathrm{Ok}(C_0)\). Fix a rank-three space \(S\), a point \(x\)
and its given chart \(L\). We write \(A=S.A\), \(B=S.B\), \(R=L.R\),
\(O'=L.O'\), \(K=O'/\mathfrak m_{O'}\), and \(d=\sum_v d_v\).
The entry map is \(\iota:A\to R\). Keep the residue embedding
\(k=O/\mathfrak m_O\hookrightarrow K\) and the normalized representation
\(r=L.\rho.r\) with its ordered characters \(\psi_{v,i}\).
The three weights are \((1,0,-1)\). Let \(E_A\) and \(E\) denote the
ideals of all entries in the four positions \(01,10,12,21\) in \(A\)
and \(R\), respectively. Let \(T_+,T_-\) denote the ideals generated
by the entries strictly below and strictly above the diagonal, respectively,
in whichever of these two rings is specified. Thus \(T_+\not\subset x\)
means that some strictly lower entry is nonzero at the point.

% BEGIN REX1
\paragraph{REX1. The sign action on the actual chart and its local flags.}
\label{claim:REX1}
Suppose \(E_A\subset x\). There is an \(O'\)-linear ring involution
\(\sigma\) on \(R\) fixing \(B\)'s image, with
\(\sigma(r_{ij}(s))=(-1)^{w_i-w_j}r_{ij}(s)\).
The ordinary box of APD21 can be chosen so that its local changes of
basis intertwine this action with diagonal sign matrices.

\begin{proof}
Twist the original normalized universal lift by the weight torus at
\(-1\). This conjugation multiplies each entry by its indicated sign,
fixes the local characters and the diagonal mark, and leaves each
continuity and residue condition intact. Universality gives the
endomorphism of the global deformation ring. Weight-zero monomials and
the local characters generate a dense subalgebra of \(B\), hence their
images and all of \(B\) are fixed. Indeed a local algebra map preserves
each maximal-ideal power and the target is separated. The twist
restricts to \(A\), generated by \(B\) and the entries. Its square
is the identity on those generators, and its difference from the
identity lies in \(E_A\). It therefore induces the same residual point
of the specified chart. Applying the chart's universal property gives
an \(O'\)-map \(\sigma:R\to R\); applying uniqueness twice shows
\(\sigma^2=1\). In \(R/E\) the quotient map and its composite with
\(\sigma\) agree on \(A\). This quotient is a nonzero complete local
coefficient algebra since \(E_A\subset x\); chart uniqueness gives
equality on \(R\). Consequently the ideal generated by the negative
eigenspace of \(\sigma\) equals \(E\): its elements vanish modulo
\(E\) since 2 is a unit, and every generator of \(E\) has negative sign.

Here is a construction over \(R\), including its nilpotents. Write
\(P=\operatorname{diag}((-1)^{w_i})\) and choose an original ordinary
basis change \(G\) at \(v\), so \(Gr(s)G^{-1}\) is upper triangular
with the ordered characters. Set \(H=GP\sigma(G^{-1})\).
Then \(H\sigma(H)=1\). At a separating local element (APD21) the
upper triangular matrix \(T=Gr(s)G^{-1}\) has diagonal values fixed
by \(\sigma\) whose pairwise differences are units, and
\(H\sigma(T)=TH\). Reading entries below the diagonal in decreasing
order of row minus column gives \(H_{ij}=0\) there: the coefficient
of \(H_{ij}\) is a unit and all other terms further below have vanished.
Choose \(D=\operatorname{diag}(\epsilon_i)\), \(\epsilon_i=\pm1\),
so \(1+\epsilon_i H_{ii}\) are units. One of the two choices works
since their sum is 2. Then \(Y=1+HD\) is upper triangular invertible
and \(H\sigma(Y)=YD\). Replacing \(G\) by \(Y^{-1}G\) yields
\(\sigma(Y^{-1}G)=D(Y^{-1}G)P\). This upper triangular change
preserves every ordinary diagonal character. At the distinguished
complex conjugation the box has matrix part one; using the two cosets
of the sign kernel extends the entry grading to its full signed group
representation. The underlying continuous representation is unchanged.
\end{proof}
% END REX1

% BEGIN REX2
\paragraph{REX2. Minimal variables and relation spaces carry the sign.}
\label{claim:REX2}
There is an equivariant surjection \(p:T=O'[[X_1,\ldots,X_a,
Y_1,\ldots,Y_g]]\to R\), with \(X\) even, \(Y\) odd. Every element
of its kernel \(J\) has all degree-one coefficients in \(\mathfrak m_{O'}\).
Minimal relations can be chosen separately in each parity, and
\(E=(pY_1,\ldots,pY_g)\).

\begin{proof}
Use the projectors \((1\pm\sigma)/2\) to lift homogeneous bases of
\(\mathfrak m_R/(\mathfrak m_{O'}R+\mathfrak m_R^2)\) to eigenvectors
in \(\mathfrak m_R\). Evaluate power series at these elements.
They and the coefficient uniformizer generate \(\mathfrak m_R\) by
Nakayama. Successive approximation in its powers lifts each element
to an element of \(T\): choose the residue constant, then homogeneous
monomials in the variables and uniformizer for the errors of orders
\(1,2,\ldots\). Completeness gives a lift and separatedness its asserted
image. Evaluation intertwines the signs on generators and hence on all
series. Averaging a lift shows it maps each parity space onto the
corresponding eigenspace in \(R\). For \(f\in J\) the constant
coefficient is in \(\mathfrak m_{O'}\) by reduction to \(K\).
In the displayed relative cotangent quotient its degree-one terms give
a dependence among the chosen basis vectors. Every coefficient therefore
vanishes in \(K\).

The finite \(K\)-space \(J/\mathfrak m_T J\) also has homogeneous
basis lifts \(f_u\in J\) by the same projectors; Nakayama makes these
generators of \(J\). Their classes are independent in each parity,
also if a parity is absent. An odd monomial contains at least one
\(Y_j\). Sorting monomials by the first such index writes each odd
series as a finite sum \(\sum Y_j h_j\). REX1 now identifies \(E\)
as claimed.
\end{proof}
% END REX2

% BEGIN REX3
\paragraph{REX3. The two obstruction counts on the original ordinary cone.}
\label{claim:REX3}
Let \(W_j\subset M_3(K)\) be the subspace supported on entries of
parity \(j\in\{0,1\}\), and \(b_{v,j}\) its intersection with the
residual ordinary Borel. Use the signed residual adjoint action, and
write \(h_{i,j}\) for the dimension of the original strict ordinary
cone group in degree \(i\) with these data. The number \(r_j\) of
minimal relations of parity \(j\) satisfies \(r_j\le h_{2,j}\).

\begin{proof}
In any equivariant small surjection \(C'\twoheadrightarrow R\) write
\(I\) for its kernel, annihilated by \(\mathfrak m_{C'}\), and \(s'\)
for its involution. Take the continuous set lifts of APD14: global
matrix parts \(r_s'\), local basis changes \(g'_v\), and triangular
matrices \(t'_{v,s}\). In a matrix with row and column signs
\(\alpha_i,\beta_j\in\{\pm1\}\) replace entry \(u_{ij}\) by
\((u_{ij}+\alpha_i\beta_j s'(u_{ij}))/2\). REX1 ensures each entry
still has the prescribed image. The determinant has invertible residue,
zeros below the diagonal remain zero, and the operation respects
continuity modulo every power. Matrix multiplication contracts the
matching internal signs, and taking inverses exchanges row and column
signs. The polarized transpose-inverse action preserves the marking
signs since \(w_{2-i}\equiv w_i\pmod2\).

Form the two defects, with \(\theta_s\) the signed automorphism,
\[
 r'_s\theta_s(r'_t)(r'_{st})^{-1}-1,\qquad
 r'_{\mathrm{loc}_v s}\big((g'_v)^{-1}t'_{v,s}g'_v\big)^{-1}-1.
\]
Their entries lie in \(I\) and have the global position signs.
Evaluate them entrywise by \(u:I\to K\) of parity \(j\), i.e.,
zero on \(I_{1-j}\). The results are \(W_j\)-valued and locally
constant by APD15. Associativity gives its global two-cocycle equation;
the difference of its restriction and the local coboundary belongs
to the Borel by multiplication of the local triangular lifts. It lies
in \(b_{v,j}\) because both terms have parity \(j\). This gives a
linear map \(I_j^\vee\to H^2_{C,j}\). A zero class supplies the
global correction cochain and the local change-of-basis corrections
in the same parity. Forgetting parity gives exactly the correcting
data of APD16--18. For every permitted quotient of \(C'\) killing
\(\ker u\subset I\), those arguments correct multiplication, restore
the ordinary box and mark, and give a section to \(R\). The section
here is an ordinary coefficient-algebra section, as required.

Apply this construction to \(C'=T/\mathfrak m_TJ\), whose kernel over
\(R\) is \(I=J/\mathfrak m_TJ\). A nonzero \(u\in I_j^\vee\) cannot
have zero obstruction. Otherwise quotienting \(C'\) by \(\ker u\)
would give a split extension with one-dimensional kernel. Indeed
\(\ker u\) is an ideal since the ring acts on \(I\) through \(K\).
To see the contradiction, write \(q:T\to C''\) for that quotient map
and \(s:R\to C''\) for a proposed section. For a relation \(f\in J\),
the images under \(q\) and \(sp\) differ on variables by elements of
the small kernel, killed by \(\mathfrak m_{C''}\). The constant terms
agree; linear terms have coefficients in \(\mathfrak m_{O'}\), and
each term of degree at least two acquires at least one further factor
from the maximal ideal. Thus \(q(f)=sp(f)=0\), using continuity
for infinite sums. This kills the whole kernel of \(C''\to R\),
a contradiction. The obstruction map is injective on \(I_j^\vee\).
REX2 identifies its dimension as \(r_j\).
\end{proof}
% END REX3

% BEGIN REX4
\paragraph{REX4. The rank-three exceptional presentation.}
\label{claim:REX4}
For every chart above with \(E_A\subset x\), there are \(b\) even
relations \(F_i\) and \(g\) odd relations \(G_j\) such that
\[
 R\simeq T/(F_1,\ldots,F_b,G_1,\ldots,G_g),
 \qquad a=b+3d+1.
\]
This gives every field of \(L.\mathrm{EPres}\), including the
coefficient-ring map, original entry lifts, and the exceptional ideal.

\begin{proof}
The adjoint respects the two matrix blocks for parity. Its open
kernel, local Borel stability and vanishing of the cyclotomic dual
from APD22--23 pass to each summand. For the last assertion, extend
a functional by zero on the complementary stable summand. The maps
cutting out the two block patterns intertwine the adjoint. In the local
basis of REX1 each basis vector has a sign. The involution acts as
the identity on the residue field since it fixes \(O'\), which maps
onto that field. Reducing \(\sigma(G')=DG'P\) modulo the maximal
ideal therefore exhibits \(D\) and \(P\) as conjugates over \(K\).
Their sign multiplicities are two and one. Of the upper triangular matrix units, three
diagonals and one off-diagonal pair of equal signs are even; the
other two are odd. This gives local dimensions \(4,2\).
At a real involution let \(U\) be its matrix part and \(J_3\) the
antidiagonal identity. The signed involution identities give
\(UJ_3=J_3U^t\); invariance of \(X\) gives
\(XU=-UJ_3X^tJ_3\). Thus \(XUJ_3\) is skew-symmetric. Since
\(UJ_3\) is even and invertible, this embeds the two invariant
spaces in the even and odd skew spaces, of dimensions one and two,
respectively (2 is invertible). The literal Euler inequality in
\(F.\mathrm{Ok}.ec\), using \(\#\mathrm{Inf}=d\), gives finite
\(h_{2,j}\) and
\[
 h_{2,0}+h_{0,0}+3d\le h_{1,0},\qquad
 h_{2,1}+h_{0,1}\le h_{1,1}.
\]

For clarity, retain the strict local data in the tangent calculation.
A strict parity cocycle and its local primitives can first be lifted
in the full adjoint by APD19. If \(D:R\to K\) is the resulting
tangent derivation, replace it by
\(D_j=(D+(-1)^j D\circ\sigma)/2\), and project the accompanying
coboundary matrix onto \(W_j\). The normalized cocycle of \(D_j\)
is the projection of that of \(D\): on an entry of the global
matrix the formula selects its position parity, and right
multiplication by the inverse residual matrix preserves parity.
The full local rigidity of APD22 descends since
\(b_{v,j}=W_j\cap b_v\). APD6 therefore gives
\(h_{1,j}\le\dim T_j\) for these tangent spaces. Its extra-direction
clause applies in the even part: the nonzero marked torus direction
\(\operatorname{diag}(1,0,-1)\) of APD20 is even and its cocycle is
unchanged by projection. Hence \(h_{1,0}+1\le\dim T_0+h_{0,0}\).
A tangent derivation is determined by the relative cotangent space
and kills variables of opposite sign. Thus \(\dim T_0\le a\),
\(\dim T_1\le g\). Combining with REX3 yields
\[
 r_0+3d+1\le a,\qquad r_1\le g .
\]
Pad the minimal lists with zeros to lengths \(b=a-3d-1\) and \(g\).
Their common kernel ideal is still \(J\), and its linear coefficients
satisfy REX2's minimality. For each entry of \(r\), averaging any
series lift puts it in the required weight parity; for each element
of the original \(B\) use the even projector. REX2 gives the remaining
ideal equality. Neither bigness nor a nontriangularity assumption
entered the construction.
\end{proof}
% END REX4

% BEGIN REX5
\paragraph{REX5. An odd tangent on the exceptional, lower-triangular chart.}
\label{claim:REX5}
Assume also \(T_+\not\subset x\). Then \(g>0\) in REX4.
\begin{proof}
In the residue the matrices are lower triangular by the second branch
of APD21 (a lower entry survives); the exceptional hypothesis kills
all four odd positions. Thus, on \(G_0\) the diagonal is
\((\Psi,1,\Psi^{-1})\) with only the \(20\) position free.
In \(W_1\) the subspace \(U=\langle e_{10},e_{21}\rangle\) is stable.
The \(20\) matrix unit commutes with both generators, so \(G_0\)
acts there by \(\Psi^{-1}\); on \(W_1/U\) it acts by \(\Psi\).
At the distinguished conjugation \(c\), the adjoint is
\(X\mapsto-J_3X^tJ_3\). Thus the sum of the two lower generators
is a \(\Psi^{-1}\delta\) line and their difference supplies a
stable complement, including in characteristic three.

At a local separating element the first two spaces of the ordered
flag are the line \(K(e_0+\lambda e_2)\) and that line plus
\(Ke_1\): these are the successive eigenspaces with characters
\(\Psi,1\) of the displayed lower-triangular matrix with distinct
eigenvalues. If \(u e_{10}+v e_{21}\) preserves the flag, applying
it to \(e_0+\lambda e_2\) gives \(u=0\), and then applying it to
\(e_1\) gives \(v=0\). Thus \(U\cap b_{v,1}=0\) at every place.

Include the strict cone for this sum line, with zero local subspaces,
in the odd cone. The induced map on \(H^1\) is injective. Indeed,
if a pair in the line is a boundary in \(W_1\), write its global
component as \((\operatorname{Ad}-1)X\) and its local primitives
as \(X\) modulo \(b_{v,1}\). The image of \(X\) in \(W_1/U\)
is invariant, hence zero since \(\Psi(g_0)\ne1\). The primitives,
already in the line, now equal \(X\) at each place because
\(U\cap b_{v,1}=0\). Projecting \(X\) in \(U\) onto the sum line
gives a strict boundary in the original line (this also handles
an empty set of places). The original \(F.\mathrm{Ok}.f4\) supplies
a nonzero class over \(k\); it survives over \(K\) because a
\(k\)-linear retraction \(K\to k\) of the specified embedding
commutes with these scalar character actions and the local conditions.
Hence \(0<h_{1,1}\le\dim T_1\le g\) by REX4.

This proves positivity for our presentation. More generally if an
arbitrary original \(L.\mathrm{EPres}\) had no odd variables, its
field identifying \(E\) would force \(E=0\). But in our presentation
the images of the \(Y_j\) are members of a relative cotangent basis,
and at least one exists. Thus \(E\ne0\), which is intrinsic to the
chart. This gives positivity for every original presentation on this
exceptional branch.
\end{proof}
% END REX5

% BEGIN REX6
\paragraph{REX6. Determinant incidence and the component graph.}
\label{claim:REX6}
On the chart of REX5 take the REX4 presentation. Each odd relation
can be written \(G_i=\sum_jM_{ij}Y_j\) by assigning its monomials to
their first odd variable. The resulting \(M_{ij}\) are even with
constant coefficient in \(\mathfrak m_{O'}\), by minimality. Put
\(\Delta=p(\det M)\) and \(I^\sharp=(\Delta)\subset R\). Then
\(I^\sharp E=0\), so \(I^\sharp\subset P\) for every prime
\(P\) outside \(V(E)\). Every minimal component of \(R/I^\sharp\)
has dimension at least \(3d+2\); any two can be joined by a chain
of such components in which successive intersections have dimension
at least \(3d+1\).
\begin{proof}
The adjugate identity applied to \(MY=G\) gives \(\Delta p(Y_j)=0\)
in \(R\), proving the first assertion.

Here is the dimension calculation behind the chain. Start in the
polynomial ring over the complete regular local domain \(T\) with
a generic \(g\)-by-\(g\) matrix \(Z\) and vector \(v\), and take
\(J_0=(Zv,\det Z)\). On \(v_i\ne0\) the \(g\) equations \(Zv=0\)
solve the \(i\)th column freely in terms of the others; the determinant
then vanishes. These opens are integral and have nonempty pairwise
intersections, dense in both members. A point with \(v=0\) and
\(\det Z=0\) over its residue field lies in their union's closure:
choose a nonzero vector \(u\) in the kernel and specialize the family
\((Z,tu)\) at \(t=0\). This works at every prime of the base. Thus
\(H=\sqrt{J_0}\) is a homogeneous prime. Its height is \(g\), as one
can localize at \(v_1\notin H\), where the eliminated column is a
family of \(g\) independent polynomial variables after an invertible
linear change.

Localize at the vertex maximal ideal \((\mathfrak m_T,Z,v)\).
Its ambient regular local ring has dimension \(\dim T+g^2+g\).
It is catenary (it is smooth over the excellent complete regular
local ring \(T\)); the local quotient by \(H\) has dimension
\(\dim T+g^2\). Completion preserves local dimension and preserves
the primality of this particular \(H\): give \(Z,v\) degree one
and \(T\) degree zero. Each homogeneous quotient piece is a finite
complete \(T\)-module. The vertex completion of the graded quotient
is the product of its pieces, since at truncation order \(N\)
the degree \(i<N\) coefficient is taken modulo
\(\mathfrak m_T^{N-i}\) and all higher degrees disappear. First
nonzero homogeneous terms therefore multiply to a nonzero term.
A power of \(H\) is already in \(J_0\) by Noetherianity. Consequently
the completed quotient by \(J_0\) has prime nilradical and dimension
\(\dim T+g^2\).

In that complete local ring impose \(Z_{ij}-M_{ij},v_j-Y_j,F_i\).
They belong to the maximal ideal; evaluation of power series is
defined as all substitution values lie in \(\mathfrak m_T\).
Its kernel on the ambient power series ring is generated by the
listed differences of variables (translate the variables by these
values and then set them to zero). The resulting ring is precisely
\(R/(\Delta)\). There are \(g^2+g+b\) cuts, leaving
\[
 (1+a+g+g^2)-(g^2+g+b)=3d+2.
\]
Apply the prime-nilradical connectedness theorem ACC2 to these cuts.
The quotient correspondence of ACC3 transports both dimension
bounds to the claimed minimal primes of \(I^\sharp\).
\end{proof}
% END REX6

% BEGIN REX7
\paragraph{REX7. The actual thick representation along the determinant.}
\label{claim:REX7}
For any prime \(P\supset(E,\Delta)\) that lies outside both triangular
loci, the literal \(\mathrm{Thick}(L.\rho,P)\) holds.
\begin{proof}
Let \(Q=R/P\), write bars for images there, and choose
\(0\ne\nu\in Q^g\) with \(\overline M\nu=0\). A kernel vector
exists over the fraction field, and its finitely many coordinates
have a common denominator. The derivation
\[
 D_0 f=\sum_j\overline{p(\partial f/\partial Y_j)}\,\nu_j
\]
vanishes on every even series because each derivative is odd and
every \(\bar Y_j\) vanishes. It kills \(G_i\) by Leibniz and the
kernel-vector identity. It descends to \(D:R\to Q\), with
\(D(pY_j)=\nu_j\). This holds for the complete power series, using
formal differentiation followed by the given quotient maps.
It kills all even entries and all the local characters by REX4.
Some odd entry has nonzero derivative: otherwise Leibniz and
\(\bar E=0\) would make \(D\) vanish on all of \(E\), contrary to
the nonzero coordinate of \(\nu\).

Define
\[
 \widehat r(s)=
 \begin{pmatrix}
 \bar r_{00}(s)&D r_{01}(s)&\bar r_{02}(s)\\
 0&\bar r_{11}(s)&0\\
 \bar r_{20}(s)&D r_{21}(s)&\bar r_{22}(s)
 \end{pmatrix}.
\]
For matrices whose odd entries reduce to zero and whose even entries
have zero derivative, this operation is multiplicative. In the
outer-to-middle column the two terms in Leibniz's rule are exactly
the product of the outer block with the new column and the product
of the new column with the middle scalar. In even positions the
products through an odd position vanish after reduction. This also
preserves the identity. Hence \(\widehat r\) is a representation.
The same calculation commutes with subtracting scalar matrices from
the ordinary factors, since \(D\psi=0\). Applying it to their ordered
product proves the required ordinary relation. The original polarized
identity \(r(s^*)=J_3 r(s)^t J_3\) reverses the factor order and
the character order; applying the construction to that reversed
identity gives the second required ordinary relation. Explicitly, for
\(s_0,s_1,s_2\in G_v\), the two ordered products are
\[
 \prod_{i=0}^{2}
 (\widehat r(\mathrm{loc}_{0,v}s_i)-\bar\psi_{v,i}(s_i)1)=0,
 \qquad
 \prod_{i=0}^{2}
 (\widehat r((\mathrm{loc}_{0,v}s_i)^*)-\bar\psi_{v,2-i}(s_i)1)=0,
\]
with factors ordered from left to right by increasing \(i\). The derivatives
and the reduced entries are adically continuous, since by Leibniz
\(D(\mathfrak m_R^{N+1})\subseteq(\mathfrak m_R Q)^N\) for every
\(N\). This proves the entrywise continuity for the original topology.

Suppose that after an injective \(Q\to K'\) into a field, the middle
character has an invariant lift \((z_0,1,z_2)\). Subtract its
eigenvector equation at \(g_0^*\) from the equation at \(g_0\).
The marked difference is diagonal, with entries \(\zeta,0,-\zeta\)
and \(\zeta\) a unit since its residue is
\(\Psi(g_0)-\Psi(g_0)^{-1}\ne0\). Its off-diagonal derivatives in
that difference are zero. Thus \(z_0=z_2=0\). For every \(s\), the
eigenvector equation now kills the derivatives in the \(01,21\)
entries. Polarization at \(s^*\) kills those in the \(12,10\) entries
as well. Injectivity gives the same vanishing over \(Q\), a
contradiction. Together with the two zero middle-row entries and
the equality of all even entries, this proves the full thick witness.
\end{proof}
% END REX7

% BEGIN REX8
\paragraph{REX8. The rank-three strict triangular fibre has at most \(2C_0\) parameters.}
\label{claim:REX8}
Put \(Q=A/(\mathfrak m_B A+T_-)\). Then \(Q\) is generated over
\(k\) by at most \(2C_0\) elements. For every original chart with
\(T_+\not\subset x\), this implies \(\dim R/\mathfrak m_B R\le2C_0\).
\begin{proof}
RTC10 gives the coefficient field action, the assigned diagonal and
local character values, and locally constant entries in \(Q\).
The case \(Q=0\) is immediate. For \(s\in G_0\), put \(p_s=\Psi(s)\).
The lower triangular residual shape, including polarization, is
\[
 r(s)=\begin{pmatrix}p_s&0&0\\ y(s)&1&0\\ z(s)&y(s^*)&p_s^{-1}\end{pmatrix}.
\]
Consequently
\[
 \begin{split}
 y(st)&=p_t y(s)+y(t),\qquad y(g_0)=y(g_0^*),\\
 z(st)&=p_t z(s)+y(s^*)y(t)+p_s^{-1}z(t),\qquad z(s^*)=z(s).
 \end{split}
\]
At each place write the lower triangular matrix of a separating element
in a lower unitriangular eigenbasis. This basis exists over the ring
\(Q\) itself: solve for the entries below each eigenvector's leading
1, dividing only by the differences of the three distinct elements
\(p_t,1,p_t^{-1}\) of \(k\). In this basis all local matrices are
still lower triangular. They are also upper triangular with ordered
diagonals by the projector computation of ACP3: placing \(s\) in
position \(i\) of the ordinary product, with the separator in all
other positions, kills the \(ij\)-entry for \(i>j\) and identifies
the diagonal. Here the eigenbasis was constructed explicitly, so the
projector cancellation works over this possibly nonlocal \(Q\).
Writing the lower entries of the basis matrix as
\(\alpha_v,\gamma_v,\beta_v\) in positions \(10,20,21\), conjugation
of \(\operatorname{diag}(p,1,p^{-1})\) gives at the local element \(s\)
\[
 y(s)=(p-1)\alpha_v,\quad y(s^*)=(1-p^{-1})\beta_v,\quad
 z(s)=(p-p^{-1})\gamma_v-(1-p^{-1})\alpha_v\beta_v.
\]

The scalar-valued functions satisfying the displayed global and local
linear conditions on \(y\) form a space of dimension at most \(C_0\).
Indeed \(y/p\) is a cocycle on \(G_0\) for \(\Psi^{-1}\).
On the two local embeddings into \(G_0\), the strict primitives are
\(-\alpha_v\) and \(\beta_v\), respectively. For the conjugate embedding
use \(csc^{-1}=(s^{-1})^*\) and \(\Psi(csc^{-1})=p\).
The primitives are determined by a separator. A zero class gives
\(y(s)=h(1-p_s)\); the equality at \(g_0,g_0^*\) forces \(h=0\)
since \(p_{g_0}\ne p_{g_0}^{-1}\). This embeds the space in the
literal group bounded by \(F.\mathrm{Ok}.b1\).

Locally constant functions on the compact source have finite image.
Choose a basis of the finite-dimensional span of the values of \(y\).
Its coordinate functions are linearly independent (otherwise the values
lie in a proper subspace), and satisfy the scalar conditions just proved.
The local parameters may be taken in the same span by evaluating at a
separator and dividing by its nonzero scalar. Thus \(y\)'s values
span at most \(C_0\) dimensions. Let \(U\subset Q\) be the
\(k\)-subalgebra generated by a basis of these values.
Both \(\alpha_v\) and \(\beta_v\) are in \(U\). Pass to the
vector-space quotient \(Q/U\). The quadratic term \(y(s^*)y(t)\)
and the local \(\alpha_v\beta_v\) term disappear. Consequently
the image of \(z\) satisfies exactly the three strict linear
conditions for \(F.\mathrm{strS}(Q/U)\) in ACC7, including the
star identity. ACC7 and RTC13 bound the span of its values by
\(C_0\). Lift a basis of those values to \(Q\).
Together with the generators of \(U\) these lifts generate all
three lower entries. RTC10 shows that they generate \(Q\), proving
the assertion also for a nonreduced fibre algebra.

For the local inequality take any prime \(P\supset\mathfrak m_B R\).
Some lower entry of weight \(-a<0\) is a unit in \(R\). Given an
upper entry of weight \(b>0\), its \(a\)th power times the lower
entry's \(b\)th power is balanced and belongs to \(\mathfrak m_B\).
Hence that upper entry vanishes modulo \(P\). The contraction
\(\eta=P\cap A\) contains \(T_-+\mathfrak m_BA\), so \(A/\eta\)
has transcendence degree at most \(2C_0\) over \(B/\mathfrak m_B=k\)
by the generator count. ACC5 on this very chart bounds \(\dim R/P\)
by that transcendence degree plus \(\dim B/\mathfrak m_B=0\).
Taking minimal primes proves the second assertion. Only the prime
quotients were used to kill the upper entries of this completed fibre.
\end{proof}
% END REX8

% BEGIN REX9
\paragraph{REX9. The opposite triangular locus is too small.}
\label{claim:REX9}
For \(T_+\not\subset x\) in rank three one has
\[
 T_+R=R,\qquad \dim R/T_-R\le2d+3C_0+2 .
\]
In particular \(F.\mathrm{Big}(C_0)\) gives \(\dim R/T_-R<3d\).
\begin{proof}
Take the original ACC9 detectors \(s_a\) (\(1+C_0+d\) of them),
\(s'_a\) (\(d\) of them), and their common positive exponent.
Cut \(R/T_-R\) by a uniformizer of \(O\), the \(r_{00}(s_a)-1\),
and the \(r_{11}(s'_a)-1\). These are \(2d+C_0+2\) elements of
the maximal ideal by their prescribed residues. In any prime
quotient \(D\) after these cuts we have a characteristic-\(\ell\)
local domain with its embedded \(k\). The three lower triangular
diagonal entries are characters. The first divided by \(\Psi\) is
residually trivial, kills all \(s_a\), and hence is trivial by the
first detector of ACC9 and the finite-quotient argument of ACC10.
For the middle diagonal character, polarization gives
\(r_{11}(g^*)=r_{11}(g)\). It kills the \(s'_a\); the second detector
and ACC10 now give \(r_{11}=1\). Polarization of the outer entries
forces the last character to equal \(\Psi^{-1}\).

The given three local characters agree with these diagonal restrictions
in the specified order. At a separator \(t\) the characteristic polynomial
has the three distinct diagonal values, and over the domain a root must
be one of them. The specified residues fix the matching at \(t\).
For arbitrary \(s\), apply the trace identity of ACP4 to
\(s,ts,t^2s\). The resulting Vandermonde system in the differences
of the two lists of character values has unit determinant, so each
difference is zero. A balanced monomial involving a lower entry
must involve an upper entry and vanishes in \(D\); the other balanced
monomials and the local characters have their assigned residual values.
The generating subalgebra of \(B\) therefore maps into the residual
constants. Adic density and continuity, as in ACC11, imply that
\(\mathfrak m_B R\) maps to zero in every such \(D\). Consequently
the quotient after the cuts has dimension at most that of
\(R/\mathfrak m_B R\), bounded by REX8. The local height bound of
ACC1 gives \(2d+C_0+2+2C_0\) before the cuts. Finally one lower
entry is a unit, so \(T_+R=R\). The bigness inequality
\(d>3C_0+7\) yields the claimed strict bound. This proves the
original rank-three branch of \(\mathrm{l51}\), with five units
of slack in its first requested bound.
\end{proof}
% END REX9

% BEGIN REX10
\paragraph{REX10. Full rank-three sharpness.}
\label{claim:REX10}
Assume \(F.\mathrm{Ok}(C_0)\), \(F.\mathrm{Big}(C_0)\),
\(S\) of rank three, and \(T_+\not\subset x\) at the given chart
\(L\). There is an ideal \(I^\sharp\subset R\) with the full
original \(L.\mathrm{Sharp}\) record.
\begin{proof}
If \(E_A\subset x\), use \((\Delta)\) from REX6. Its containment
in every prime off \(E\) and the required dimension and incidence
statements are already proved there. For a minimal prime \(P\)
over it, \(\dim R/P\ge3d+2\). It cannot contain \(T_-R\)
by REX9, and cannot contain \(T_+R=R\). It is consequently type I
if it omits \(E\), and otherwise satisfies all premises of REX7,
which gives the original thick witness. If \(E_A\not\subset x\),
one of the odd generators is a unit and \(E=R\). Use \(I^\sharp=0\).
The ordinary presentation and ACC4 give the component dimensions
and chain at \(3d+1\). REX9 again excludes both triangular ideals
from every minimal prime, so every such prime is type I. The
zero ideal has the required universal containment off \(E\).
These are the four fields of \(\mathrm{Sharp}\) on the supplied
chart. The distinct rank-two support ideal and the rank-three
corner ideal used further in propagation have not entered its proof.
\end{proof}
% END REX10
